% ============================================================================
%  Rapport Technique — Système Intelligent d'Appariement ACPE
%  Hackathon IndabaX Congo 2026 — Équipe S2M (Stanislas Mboungou Mangoubi)
%
%  Source LaTeX reconstruite et mise à jour (juillet 2026) :
%   - §3.1 : calibration logistique rendue continue en 0 (correction de la fuite logistique)
%   - §3.2 : Table 2 remplacée par le benchmark empirique TF-IDF vs BERT vs Hybrides
%   - §5   : analyse scientifique de la supériorité du TF-IDF sur les embeddings BERT
%
%  Compilation :  pdflatex Rapport_INDABA_CG_2026.tex  (x2 pour les références)
% ============================================================================
\documentclass[11pt,a4paper]{article}

\usepackage[utf8]{inputenc}
\usepackage[T1]{fontenc}
\usepackage{lmodern}
\usepackage[french]{babel}
\usepackage[margin=2.5cm]{geometry}
\usepackage{amsmath,amssymb}
\usepackage{booktabs}
\usepackage{array}
\usepackage{xcolor}
\usepackage{fancyhdr}
\usepackage{caption}
\usepackage[hidelinks]{hyperref}

% En-tête / pied de page courants (comme dans le PDF d'origine).
\pagestyle{fancy}
\fancyhf{}
\lhead{\small Système intelligent d'appariement ACPE}
\rhead{\small Hackathon IndabaX Congo 2026}
\lfoot{\small Équipe S2M~-- Stanislas Mboungou Mangoubi}
\rfoot{\small \thepage}
\renewcommand{\headrulewidth}{0.4pt}

\newcommand{\raw}{\text{Score}_{\text{raw}}}

\title{%
  \vspace{-1.2cm}
  {\normalsize \textsc{Agence Congolaise pour l'Emploi}}\\[0.6em]
  \textbf{Rapport Technique : Système Intelligent d'Appariement\\
  Demandeurs et Offreurs d'Emploi}\\[0.3em]
  {\large Solution IA Explicable et Opérationnelle pour l'ACPE}}
\author{Stanislas Mboungou Mangoubi, MSc. Actuariat\\
  \small Email : \href{mailto:stanislasaigle@gmail.com}{stanislasaigle@gmail.com}\\
  \small Hackathon IndabaX Congo 2026}
\date{Juillet 2026}

\begin{document}
\maketitle
\thispagestyle{fancy}

% ---------------------------------------------------------------------------
\section*{Résumé Exécutif}
Ce rapport présente la conception technique du moteur d'appariement intelligent développé
par l'équipe \textbf{S2M} pour l'Agence Congolaise pour l'Emploi (ACPE). Face à un volume
critique de données (\textbf{41\,285} demandeurs pour \textbf{2\,535} offres), notre solution
implémente une approche hybride unissant la \textbf{similarité textuelle vectorielle (TF-IDF
Cosinus)} à des règles métiers déterministes. Conforme aux exigences du jury, ce système se
distingue par son \textbf{explicabilité totale (XAI)} terme à terme, son équité algorithmique
native, sa capacité d'analyse d'écarts de compétences (\emph{skill gap}) et un tableau de bord
d'aide à la décision entièrement interactif et filtrable. Le choix du TF-IDF n'est pas un choix
par défaut : il est \textbf{validé empiriquement} (\S\ref{sec:benchmark}) contre une alternative
à base d'embeddings BERT, qu'il surpasse nettement sur nos données.

% ---------------------------------------------------------------------------
\section{Contexte et Problématique}
L'Agence Congolaise pour l'Emploi (ACPE) occupe une position stratégique dans la régulation
du marché du travail en République du Congo. Cependant, l'identification des meilleures
correspondances entre les profils des demandeurs et les exigences des entreprises repose
historiquement sur un traitement essentiellement manuel et chronophage. À l'ère de la science
des données, l'absence d'un outil d'aide à la décision automatisé limite la réactivité des
conseillers et la valorisation des informations collectées.

L'enjeu de ce projet est de concevoir un prototype logiciel capable de calculer un score de
compatibilité transparent, de générer des recommandations optimales de type Top-K (Top-5 et
Top-10), et de fournir une interface opérationnelle immédiate sans nécessiter d'infrastructures
de calcul lourdes ou coûteuses.

% ---------------------------------------------------------------------------
\section{Pipeline de Données et Préparation}
L'intégration et la réconciliation des quatre bases de données fournies s'effectuent au sein
d'un pipeline unifié (\texttt{src/data\_loader.py}). Le tableau~\ref{tab:donnees} résume les
volumes et la structure après nettoyage complet.

\begin{table}[h]
\centering
\caption{Structure et volumétrie des bases de données réconciliées}
\label{tab:donnees}
\small
\begin{tabular}{@{}llrl@{}}
\toprule
\textbf{Fichier Source} & \textbf{Identifiant Clef} & \textbf{Volume} & \textbf{Rôle dans le Pipeline} \\
\midrule
\texttt{Demandeurs .xlsx}                 & Matricule       & 41\,285 & Profils textuels des demandeurs \\
\texttt{Offres\_ACPE.xlsx}                & Référence offre & 2\,535  & Table de référence des offres \\
\texttt{Offres\_ACPE\_Extensions.xlsx}    & Référence       & 143     & Enrichissement textuel libre \\
\texttt{Appariement\_...\_Offres.xlsx}    & id\_demandeur   & 41\,285 & Vérité terrain (\emph{Ground Truth}) \\
\bottomrule
\end{tabular}
\end{table}

\subsection{Décisions de Nettoyage et Alignement Structurant}
\begin{itemize}
  \item \textbf{Normalisation syntaxique :} les en-têtes contenant des espaces parasites ou des
  accents (ex.\ \texttt{'offre\_pertinente~'}) ont été systématiquement nettoyés pour éviter les
  échecs silencieux d'indexation.
  \item \textbf{Fusion textuelle par complétude :} une analyse de cardinalité a prouvé que les
  143 lignes du fichier d'extension ne constituaient pas de nouvelles offres mais un
  enrichissement sémantique d'offres déjà existantes dans la table principale ; elles ont donc
  été fusionnées par jointure gauche.
  \item \textbf{Traitement des sentinelles :} la chaîne « Non déclaré », omniprésente dans les
  données ($\sim 91\%$ pour le champ \emph{Secteur demandé}), est traitée explicitement comme une
  valeur manquante (\texttt{NaN}) afin de ne pas biaiser la fréquence des termes du modèle.
\end{itemize}

\subsection{Simulation Déterministe de la Position Géographique}
Le fichier initial des demandeurs ne présentant aucune variable de localisation exploitable,
nous avons généré une variable synthétique de localité candidate sans altérer les fichiers
sources. Afin d'assurer le réalisme du prototype, le tirage respecte strictement la
\textbf{distribution probabiliste réelle} des offres (\texttt{Lieu}), concentrant ainsi les
candidats simulés sur les bassins d'emploi effectifs (Brazzaville, Pointe-Noire, etc.). La
reproductibilité est garantie par un hachage déterministe basé sur le \texttt{Matricule} du
candidat. Cette variable simulée demeure un artefact analytique (carte du tableau de bord,
exploration) : \textbf{elle n'entre jamais dans le calcul du score} d'appariement.

% ---------------------------------------------------------------------------
\section{Architecture du Moteur d'Appariement}
Pour répondre simultanément aux critères de performance algorithmique (30\% de la note) et
d'explicabilité, nous avons implémenté un modèle hybride combinant le traitement automatique du
langage naturel (NLP) et des contraintes métiers.

\subsection{Modélisation Mathématique}
\label{sec:math}
Le score brut d'appariement entre un candidat $c$ et une offre $o$ est défini par l'équation
linéaire suivante :
\begin{equation}
\raw(c, o) = w_t \cdot \mathrm{Sim}_{\cos}(v_c, v_o) + w_m \cdot \mathrm{Recouvrement}(M_c, I_o)
\end{equation}
Où :
\begin{itemize}
  \item $v_c$ et $v_o$ représentent les vecteurs \textbf{TF-IDF} du candidat et de l'offre,
  calculés sur un espace de n-grammes de longueur 1 à 2 (\texttt{ngram\_range=(1,2)}), filtrant
  les termes apparaissant moins de deux fois (\texttt{min\_df=2}).
  \item $\mathrm{Sim}_{\cos}$ désigne la similarité cosinus mesurant l'alignement directionnel
  des profils :
  \begin{equation}
  \mathrm{Sim}_{\cos}(v_c, v_o) = \frac{v_c \cdot v_o}{\lVert v_c \rVert\,\lVert v_o \rVert}
  \end{equation}
  \item $\mathrm{Recouvrement}(M_c, I_o)$ mesure l'intersection normalisée des tokens entre le
  \emph{Métier visé} ($M_c$) et l'\emph{Intitulé de l'offre} ($I_o$), agissant comme un
  amplificateur de précision pour les correspondances exactes de l'intitulé du poste.
  \item Les poids empiriques sont fixés à $w_t = 0{,}80$ et $w_m = 0{,}20$.
\end{itemize}

Pour l'affichage utilisateur, le score brut est projeté dans l'espace $[0, 100]$ via une
fonction logistique de calibration cosmétique, strictement monotone, préservant parfaitement le
rang. \textbf{Correction de la fuite logistique.} La sigmoïde brute mappe un score nul vers une
valeur \emph{non nulle} : $100 / (1 + e^{0{,}15 \cdot 8}) \approx 23{,}1\%$. Or un score brut nul
signifie une orthogonalité totale des vecteurs (aucun token en commun), c'est-à-dire une
\emph{absence de correspondance} qui ne doit pas se lire comme $23\%$. Nous imposons donc une
condition explicite $\mathrm{Compatibilité}(0) = 0$, rendant la calibration continue en zéro tout
en restant monotone :
\begin{equation}
\mathrm{Compatibilité}(\%) =
\begin{cases}
0 & \text{si } \raw \leq 0, \\[6pt]
\dfrac{100}{1 + e^{-8\,(\raw - 0{,}15)}} & \text{si } \raw > 0.
\end{cases}
\label{eq:calib}
\end{equation}
Cette condition élimine les « faux positifs à $23\%$ » et fonde le \textbf{seuil opérationnel
d'affichage} ($45\%$) appliqué aux onglets de recherche : toute recommandation dont la
compatibilité calibrée est inférieure à ce seuil est masquée plutôt que présentée au conseiller.

\subsection{Justification Scientifique face aux Modèles Alternatifs}
\label{sec:benchmark}
Le choix de l'architecture hybride TF-IDF + Règles Métiers résulte d'un arbitrage rigoureux.
Deux familles sont écartées sur des bases théoriques : les \textbf{règles métiers pures} (trop
rigides — les taxonomies de secteurs divergent entre les fichiers, d'où une couverture médiocre),
et le \textbf{Learning-to-Rank supervisé} (risque extrême de surapprentissage : la vérité terrain
ne fournit que 3 positifs par candidat, sans négatifs explicites).

L'alternative la plus sérieuse — les \textbf{embeddings sémantiques profonds (BERT / SBERT)} — a,
elle, été \textbf{évaluée empiriquement} plutôt que rejetée sur principe. Nous avons encodé les
2\,535 offres et 41\,285 profils avec un modèle multilingue léger et adapté au CPU
(\texttt{paraphrase-multilingual-MiniLM-L12-v2}, 384 dimensions), puis comparé, sur la vérité
terrain, le TF-IDF, le BERT pur et plusieurs hybrides $\alpha \cdot \text{TF-IDF} +
(1-\alpha)\cdot\text{BERT}$. Le tableau~\ref{tab:benchmark} présente ce benchmark.

\begin{table}[h]
\centering
\caption{Benchmark empirique des moteurs sur un échantillon de 4\,000 demandeurs (vérité terrain
ACPE, CPU, embeddings précalculés). $\alpha$ = poids du TF-IDF dans l'hybride.}
\label{tab:benchmark}
\small
\begin{tabular}{@{}lrrrrrr@{}}
\toprule
\textbf{Moteur} & \textbf{P@5} & \textbf{P@10} & \textbf{R@5} & \textbf{R@10} & \textbf{NDCG@5} & \textbf{NDCG@10} \\
\midrule
\textbf{TF-IDF (déployé)}          & \textbf{0,423} & \textbf{0,256} & \textbf{0,704} & \textbf{0,852} & \textbf{0,680} & \textbf{0,748} \\
BERT (MiniLM multilingue)          & 0,185 & 0,125 & 0,308 & 0,415 & 0,279 & 0,328 \\
Hybride $\alpha=0{,}3$             & 0,319 & 0,195 & 0,532 & 0,651 & 0,532 & 0,586 \\
Hybride $\alpha=0{,}5$             & 0,354 & 0,216 & 0,590 & 0,720 & 0,595 & 0,655 \\
Hybride $\alpha=0{,}7$             & 0,374 & 0,227 & 0,623 & 0,757 & 0,622 & 0,683 \\
\bottomrule
\end{tabular}
\end{table}

Le verdict est sans appel : le TF-IDF domine sur toutes les métriques, et — signe décisif — la
performance de l'hybride \textbf{croît de façon monotone avec $\alpha$}. L'analyse scientifique de
ce résultat contre-intuitif est développée en \S\ref{sec:analyse}. Retenu : le
\textbf{modèle hybride S2M (TF-IDF + Règle)}, explicable, léger (exécution CPU en $\sim$85 s pour
41\,285 lignes) et empiriquement supérieur.

% ---------------------------------------------------------------------------
\section{Ingénierie de l'Explicabilité (XAI) et Équité}
Conformément aux orientations du jury, l'explicabilité n'est pas une surcouche mais un pilier
natif de l'algorithme :
\begin{enumerate}
  \item \textbf{Décomposition des composantes :} chaque conseiller ACPE voit la contribution
  exacte de la sémantique (80\%) et de la règle métier (20\%).
  \item \textbf{Extraction des termes déterminants :} la similarité cosinus étant une somme de
  contributions terme à terme, le module \texttt{src/explain.py} isole et affiche les mots exacts
  (ex.\ « transit », « logistique ») ayant provoqué la décision.
  \item \textbf{Équité par conception (\emph{Fairness}) :} les variables sensibles telles que
  l'âge et le genre sont explicitement exclues de la fonction de score. L'équité est totale,
  auditable et garantie par l'architecture du code.
\end{enumerate}

% ---------------------------------------------------------------------------
\section{Résultats et Métriques d'Évaluation}
\label{sec:resultats}
Le moteur a été évalué de manière exhaustive sur l'ensemble de la population de test (41\,285
demandeurs) via le script batch \texttt{scripts/generate\_recommendations.py}. Les formules
appliquées respectent les consignes officielles :
\begin{equation}
\mathrm{Precision@K} = \frac{|O_{\text{recommandées}} \cap O_{\text{pertinentes}}|}{K}, \quad
\mathrm{Recall@K} = \frac{|O_{\text{recommandées}} \cap O_{\text{pertinentes}}|}{|O_{\text{totales\_pertinentes}}|}
\end{equation}
Le tableau~\ref{tab:perf} présente les performances globales obtenues.

\begin{table}[h]
\centering
\caption{Performances globales validées face à la vérité terrain de l'ACPE (41\,285 demandeurs)}
\label{tab:perf}
\small
\begin{tabular}{@{}lrrr@{}}
\toprule
\textbf{Seuil d'Évaluation (K)} & \textbf{Precision@K} & \textbf{Recall@K} & \textbf{NDCG@K} \\
\midrule
Top-5 Recommandations  & 0,427 & 0,712 & 0,688 \\
Top-10 Recommandations & 0,257 & 0,855 & 0,753 \\
\bottomrule
\end{tabular}
\end{table}

\paragraph{Analyse des performances.} La vérité terrain limitant structurellement le nombre
d'offres pertinentes à 3 par candidat, la $\mathrm{Precision@5}$ maximale théorique est
mathématiquement plafonnée à $3/5 = 0{,}60$. Notre score de $0{,}427$ atteint ainsi
\textbf{71,2\% du maximum théorique possible}. Le $\mathrm{Recall@10}$ de $0{,}855$ garantit que
dans plus de \textbf{85\%} des cas, les conseillers trouvent l'offre idéale dans les 10 premières
options suggérées par l'IA.

\subsection{Analyse Scientifique : pourquoi le TF-IDF surpasse BERT sur nos données}
\label{sec:analyse}
Le benchmark de la Table~\ref{tab:benchmark} montre que BERT n'atteint que $0{,}185$ en
$\mathrm{Precision@5}$, soit \textbf{moins de 44\%} de la performance du TF-IDF ($0{,}423$), et que
chaque unité de poids transférée du TF-IDF vers BERT dans l'hybride \emph{dégrade} le score.
L'apport sémantique de BERT est donc, sur ce corpus, non pas neutre mais \textbf{net négatif}.
Ce résultat, contre-intuitif au premier abord, s'explique rigoureusement par trois mécanismes.

\begin{enumerate}
  \item \textbf{Pénurie de contexte sur des intitulés de 2 à 4 mots.} Les modèles de type SBERT
  sont entraînés à encoder la sémantique \emph{de phrases} via l'auto-attention sur un contexte
  étendu. Nos champs (« Agent de transit », « Statisticien ») sont des \emph{listes de mots-clés}
  quasi dépourvues de contexte : l'attention n'a rien à contextualiser et les vecteurs se
  concentrent dans un cône anisotrope étroit où le pouvoir discriminant s'effondre (phénomène
  connu de l'anisotropie des représentations, cf.\ Ethayarajh, 2019 [7]). À l'inverse, le TF-IDF
  est \emph{optimal} précisément sur des documents courts : la pondération IDF (Jones, 1972 [3])
  maximise la spécificité des tokens rares et distinctifs (« transit », « statisticien »).

  \item \textbf{Décalage de domaine (hors-distribution).} Le MiniLM multilingue est pré-entraîné
  sur des corpus web/paraphrase généralistes, \emph{non} sur la taxonomie du marché de l'emploi
  congolais. Les intitulés français spécialisés et la terminologie locale sont hors-distribution :
  l'espace dense rapproche des rôles sémantiquement voisins mais faux (\emph{semantic smearing}),
  d'où l'effondrement de la précision. Le TF-IDF, lui, est \textbf{adaptatif au corpus} : l'IDF est
  estimé directement sur les données de l'ACPE, ce qui lui confère une pertinence de domaine que
  l'embedding pré-entraîné ne possède pas.

  \item \textbf{Alignement avec la nature de la vérité terrain.} L'appariement de référence de
  l'ACPE repose massivement sur une correspondance \emph{lexicale} métier~$\leftrightarrow$~intitulé.
  Les n-grammes du TF-IDF capturent directement ce signal, tandis que la généralisation sémantique
  de BERT travaille \emph{contre} l'objectif : elle récupère des offres plausibles mais absentes de
  la vérité terrain, ce que la $\mathrm{Precision}$ et le NDCG sanctionnent. La \textbf{dégradation
  monotone} de l'hybride ($0{,}374 \to 0{,}354 \to 0{,}319$ à mesure que $\alpha$ décroît) est la
  preuve empirique directe que le signal BERT s'éloigne de la vérité terrain plutôt qu'il ne la
  complète.
\end{enumerate}

\paragraph{Corollaire décisionnel.} Même à performance égale, BERT resterait disqualifié par notre
double contrainte : \emph{boîte noire} (perte de l'explicabilité XAI exigée par le jury) et
empreinte GPU/RAM incompatible avec l'hébergement gratuit. Ici, il perd \emph{en plus} sur la
précision. Le TF-IDF n'est donc pas un choix de repli mais l'option \textbf{scientifiquement
optimale} pour ce régime de données : textes courts, domaine spécialisé, vérité terrain à
alignement lexical. Un modèle spécialisé français (type CamemBERT) pourrait réduire l'écart, mais
son ampleur ($0{,}185$ contre $0{,}423$) le rend improbable à renverser, sans résoudre les
contraintes de coût ni d'explicabilité.

% ---------------------------------------------------------------------------
\section{Fonctionnalités Bonus et Exploitation Directe}
Le prototype intègre pleinement les défis avancés requis pour transformer l'expérience
utilisateur :
\begin{itemize}
  \item \textbf{Recherche intelligente en langage naturel (Bonus 1) :} un double moteur permet des
  requêtes libres sur les offres et les candidats, avec détection automatique de la localité au
  sein de la phrase. Une garde sémantique lexicale (tolérante aux racines, ex.\ « statistique »
  $\to$ « statisticien ») isole les termes métiers des termes géographiques, empêchant qu'une
  ville ne « détourne » le classement.
  \item \textbf{Analyse des écarts de compétences (Bonus 2) :} le système extrait dynamiquement les
  exigences manquantes (ex.\ Git, Power BI) pour orienter le citoyen vers des formations ciblées.
  \item \textbf{Tableau de bord décisionnel interactif :} développé avec un design respectant la
  charte chromatique nationale (Vert, Jaune, Rouge), l'onglet Dashboard permet aux conseillers de
  filtrer dynamiquement les indicateurs (KPIs, histogrammes de compatibilité, cartographies
  géographiques) par ville et secteur d'activité en temps réel.
\end{itemize}

% ---------------------------------------------------------------------------
\section{Conclusion}
La solution portée démontre qu'il est possible d'allier une haute précision mathématique à une
transparence totale indispensable pour une agence publique. Le choix du TF-IDF, loin d'être un
compromis, est \textbf{empiriquement démontré} supérieur aux embeddings profonds sur ce type de
données. Léger, déterministe, sans dépendance matérielle lourde et immédiatement exploitable, ce
prototype pose les fondations d'une intermédiation moderne et équitable pour l'emploi en
République du Congo.

% ---------------------------------------------------------------------------
\section{Bibliographie Sélective}
\begin{thebibliography}{9}
\bibitem{salton} Salton, G., Wong, A., et Yang, C. S. (1975). \emph{A vector space model for
automatic indexing.} Communications of the ACM, 18(11), 613-620. \\
\textit{Apport clé :} fondation théorique du modèle d'espace vectoriel et de la similarité cosinus.

\bibitem{jarvelin} Järvelin, K., et Kekäläinen, J. (2002). \emph{Cumulated gain-based evaluation of
IR techniques.} ACM TOIS, 20(4), 422-446. \\
\textit{Apport clé :} invention du NDCG, métrique maîtresse de notre évaluation.

\bibitem{jones} Jones, K. S. (1972). \emph{A statistical interpretation of term specificity and its
application in retrieval.} Journal of Documentation, 28(1), 11-21. \\
\textit{Apport clé :} introduction de l'IDF pour pondérer la spécificité des mots-clés.

\bibitem{manning} Manning, C. D., Raghavan, P., et Schütze, H. (2008). \emph{Introduction to
Information Retrieval.} Cambridge University Press. \\
\textit{Apport clé :} référence pour les structures textuelles et le calcul de Precision@K.

\bibitem{lime} Ribeiro, M. T., Singh, S., et Guestrin, C. (2016). \emph{``Why should I trust
you?'': Explaining the predictions of any classifier.} Proceedings of the ACM SIGKDD. \\
\textit{Apport clé :} formalisation de l'IA explicable (XAI), justifiant le rejet des boîtes noires.

\bibitem{fairness} Barocas, S., Hardt, M., et Narayanan, A. (2019). \emph{Fairness and Machine
Learning.} fairmlbook.org. \\
\textit{Apport clé :} cadre de l'équité par conception, dictant l'exclusion de l'âge et du genre.

\bibitem{ethayarajh} Ethayarajh, K. (2019). \emph{How contextual are contextualized word
representations? Comparing the geometry of BERT, ELMo, and GPT-2 embeddings.} EMNLP-IJCNLP. \\
\textit{Apport clé :} met en évidence l'anisotropie des embeddings, expliquant l'effondrement du
pouvoir discriminant sur nos intitulés très courts (\S\ref{sec:analyse}).

\bibitem{sbert} Reimers, N., et Gurevych, I. (2019). \emph{Sentence-BERT: Sentence embeddings using
Siamese BERT-networks.} EMNLP-IJCNLP. \\
\textit{Apport clé :} modèle de référence pour les embeddings de phrases, évalué puis écarté au
profit du TF-IDF sur nos données (\S\ref{sec:benchmark}).
\end{thebibliography}

\end{document}
